\documentclass[a4paper,
fontsize=11pt,
%headings=small,
oneside,
numbers=noperiodatend,
parskip=half-,
bibliography=totoc,
final
]{scrartcl}

\usepackage[babel]{csquotes}
\usepackage{synttree}
\usepackage{graphicx}
\setkeys{Gin}{width=.4\textwidth} %default pics size

\graphicspath{{./plots/}}
\usepackage[ngerman]{babel}
\usepackage[T1]{fontenc}
%\usepackage{amsmath}
\usepackage[utf8x]{inputenc}
\usepackage [hyphens]{url}
\usepackage{booktabs} 
\usepackage[left=2.4cm,right=2.4cm,top=2.3cm,bottom=2cm,includeheadfoot]{geometry}
\usepackage[labelformat=empty]{caption} % option 'labelformat=empty]' to surpress adding "Abbildung 1:" or "Figure 1" before each caption / use parameter '\captionsetup{labelformat=empty}' instead to change this for just one caption
\usepackage{eurosym}
\usepackage{multirow}
\usepackage[ngerman]{varioref}
\setcapindent{1em}
\renewcommand{\labelitemi}{--}
\usepackage{paralist}
\usepackage{pdfpages}
\usepackage{lscape}
\usepackage{float}
\usepackage{acronym}
\usepackage{eurosym}
\usepackage{longtable,lscape}
\usepackage{mathpazo}
\usepackage[normalem]{ulem} %emphasize weiterhin kursiv
\usepackage[flushmargin,ragged]{footmisc} % left align footnote
\usepackage{ccicons} 
\setcapindent{0pt} % no indentation in captions
\usepackage{xurl} % Breaks URLs

%%%% fancy LIBREAS URL color 
\usepackage{xcolor}
\definecolor{libreas}{RGB}{112,0,0}

\usepackage{listings}

\urlstyle{same}  % don't use monospace font for urls

\usepackage[fleqn]{amsmath}

%adjust fontsize for part

\usepackage{sectsty}
\partfont{\large}

%Das BibTeX-Zeichen mit \BibTeX setzen:
\def\symbol#1{\char #1\relax}
\def\bsl{{\tt\symbol{'134}}}
\def\BibTeX{{\rm B\kern-.05em{\sc i\kern-.025em b}\kern-.08em
    T\kern-.1667em\lower.7ex\hbox{E}\kern-.125emX}}

\usepackage{fancyhdr}
\fancyhf{}
\pagestyle{fancyplain}
\fancyhead[R]{\thepage}

% make sure bookmarks are created eventough sections are not numbered!
% uncommend if sections are numbered (bookmarks created by default)
\makeatletter
\renewcommand\@seccntformat[1]{}
\makeatother

% typo setup
\clubpenalty = 10000
\widowpenalty = 10000
\displaywidowpenalty = 10000
\usepackage{accessibility}

\usepackage[colorlinks, linkcolor=black,citecolor=black, urlcolor=libreas,
breaklinks= true,bookmarks=true,bookmarksopen=true]{hyperref}
\usepackage{hyperxmp}
\usepackage{breakurl}

%meta
%meta

\fancyhead[L]{K. Brannemann, J. Kreutzer\\ %author
LIBREAS. Library Ideas, 48 (2026). % journal, issue, volume.
\href{https://doi.org/10.18452/37899}{\color{black}https://doi.org/10.18452/37899}
{}} % doi
\fancyhead[R]{\thepage} %page number
\fancyfoot[L] {\ccLogo \ccAttribution\ \href{https://creativecommons.org/licenses/by/4.0/}{\color{black}Creative Commons BY 4.0}}  %licence
\fancyfoot[R] {ISSN: 1860-7950}

\title{\LARGE{Mehr als Information: Interviews als Raum für Geschichten, Werte und Kultur}}% title
\author{Kathrin Brannemann, Josephine Kreutzer} % author

\setcounter{page}{1}

\hypersetup{%
      pdftitle={Mehr als Information: Interviews als Raum für Geschichten, Werte und Kultur},
      pdfauthor={Kathrin Brannemann, Josephine Kreutzer},
      pdfcopyright={CC BY 4.0 International},
      pdfsubject={LIBREAS. Library Ideas, 48 (2026)},
      pdfkeywords={DNB, Deutsche Nationalbibliothek, Interview, interne Kommunikation, German National Library, internal communications},
      pdflicenseurl={https://creativecommons.org/licenses/by/4.0/},
      pdfcontacturl={http://libreas.eu},
      pdfurl={https://doi.org/10.18452/37899},
      pdfdoi={10.18452/37899},
      pdflang={de},
      pdfmetalang={de}
     }



\date{}
\begin{document}

\maketitle
\thispagestyle{fancyplain} 

%abstracts

%body
Von außen betrachtet sind Organisationen oft als Strukturen sichtbar --
als Prozesse, Hierarchien, Dokumente und Zahlen. Doch die meisten
Menschen lassen sich nicht von Strukturen, sondern von Visionen und
Personen anziehen. Genau hier setzt die Kraft des Interviews an -- nicht
als bloßes Informationsmedium, sondern als ein menschlicher Blick hinter
eine Idee, eine Institution oder ein Projekt.

Interviews öffnen einen Raum, in dem Geschichten erzählt und Werte
vermittelt werden und Unternehmenskultur erlebbar wird. Die persönliche
Perspektive des/der Interviewten verleiht Glaubwürdigkeit, die schwer
durch standardisierte Texte zu erreichen ist. Sie baut Vertrauen auf,
schafft Nähe und ermöglicht eine emotionale Verbindung.

Besonders bedeutungsvoll ist dieser Ansatz in der internen
Kommunikation. Hier geht es nicht nur um Wissensmanagement, sondern um
die Verankerung von Werten, die Stärkung des Unternehmensgefühls und die
Schaffung von Identifikation. Interviews mit Führungskräften,
Projektteams oder Kolleg*innen schaffen Verständnis für
Entscheidungsprozesse.

In der Deutschen Nationalbibliothek haben wir in den vergangenen Jahren
mit verschiedenen Interviewformaten experimentiert und mittlerweile ein
festes Set entwickelt.

\section{Erfahrungen teilen: Jubilare im
Mittelpunkt}\label{erfahrungen-teilen-jubilare-im-mittelpunkt}

Ein Beispiel ist das Format \enquote{Jubilare erzählen}. Hier steht
nicht die Leistung, sondern die Lebensgeschichte im Fokus. Mit Fragen
wie: \enquote{Was hat Sie damals hierhergeführt?} und \enquote{Wie hat
sich das Arbeitsleben verändert?} teilen Kolleg*innen mit kleinen und
großen Jubiläen prägende Momente, Anekdoten und Erlebnisse. Angereichert
wird das Format durch Fragen, die einem Poesiealbum ähneln: Was macht
Ihr Büro besonders? Was ist Ihr Lieblingsessen in der Kantine? Welche
persönliche Notiz möchten Sie Kolleg*innen mitgeben?

\section{Zwischen den Regalen: Die Geschichten von
Nutzenden}\label{zwischen-den-regalen-die-geschichten-von-nutzenden}

Ein weiteres Beispiel ist die Reihe \enquote{Zwischen den Regalen}, in
der Nutzende der Deutschen Nationalbibliothek zu Wort kommen. Sie kommen
aus unterschiedlichen Bereichen, mit unterschiedlichen Zielen -- als
Forschende, Studierende, Autor*innen oder anderer Berufsbranchen. Sie
nutzen die Lesesäle in Leipzig und Frankfurt als Recherche-, Arbeits-
oder Lernort. In diesem Format kommen sie zu Wort und geben Einblicke in
ihre Arbeit, stellen Projekte vor und erzählen ihre Geschichte.
\enquote{Zwischen den Regalen} nimmt dabei häufig eine doppelte
Bedeutung ein, denn viele der Interviewten sitzen nicht nur wörtlich
zwischen Regalen, sondern so tut es auch ihr Werk -- ein bis drei
Stockwerke tiefer in der Sammlung. Auch diese werden in der Reihe
vorgestellt. So entsteht eine Verbindung zwischen dem Bestand und dessen
Nutzenden.

\section{1 von 600}\label{von-600}

Sie sind eine oder einer von 600, die die Deutsche Nationalbibliothek
mit ihrem Engagement zu Deutschlands kulturellem Gedächtnis machen. In
\enquote{1 von 600} gewähren Mitarbeitende einen Blick in ihre Arbeit
und ihr \#sinnvollesSchaffen: Mit Fragen wie: \enquote{Was darf an Ihrem
Arbeitsplatz nicht fehlen?}, \enquote{Was ist Ihr Lieblingsplatz?} und
\enquote{Wie sieht ein Arbeitstag aus?} werden Arbeitswelten intern und
extern greifbar. Kolleg*innen lernen ihre Büronachbar*innen näher
kennen, Interessierte finden Antworten auf Arbeitsabläufe und Kultur.

\section{Ein Abschied. Ein
Neuanfang.}\label{ein-abschied.-ein-neuanfang.}

Ute Schwens war Direktorin der Deutschen Nationalbibliothek in Frankfurt
am Main und ständige Vertreterin des Generaldirektors.

Sie arbeitete seit 1980 in wechselnden Verantwortlichkeiten in der
Deutschen Nationalbibliothek. 1999 hat sie die Leitung des Frankfurter
Hauses übernommen. Am 31. Januar 2026 wurde Ute Schwens in den Ruhestand
versetzt. In einem Gespräch kurz vor ihrem Austritt schaute sie zurück
auf 48 Jahre Berufserfahrung.\footnote{\url{https://blog.dnb.de/ich-habe-aufgaben-bekommen-und-geschafft-sie-zu-losen-das-starkt/}}

Kathrin Brannemann hat am 1. Februar das Amt der Direktorin der
Deutschen Nationalbibliothek in Frankfurt am Main angetreten und ist
zugleich eine der beiden ständigen Vertreter*innen des Generaldirektors
der Deutschen Nationalbibliothek, Frank Scholze. In ihrem
Antrittsinterview -- welches hier erneut publiziert wird -- teilte sie
Werte und Überzeugungen und einen Einblick auf die kommende
Aufgabe.\footnote{\url{https://blog.dnb.de/bibliotheksarbeit-wirkt-unmittelbar/}}

\section{Bibliotheksarbeit wirkt
unmittelbar}\label{bibliotheksarbeit-wirkt-unmittelbar}

\textbf{Frau Brannemann, können Sie uns zu Beginn ein wenig über sich
selbst erzählen? Welche Impulse haben Sie ursprünglich in die
Bibliothekswelt geführt?}

Meine Begeisterung für Bibliotheken hat sehr früh begonnen. Einen Teil
meiner Kindheit habe ich in den USA verbracht. Bereits in der
Grundschule hatte ich dort einen kleinen Bibliotheksdienst und als ich
meinen ersten eigenen Bibliotheksausweis bekommen habe, war das ein
prägendes Erlebnis. Ich war fasziniert davon, wie Wissen gesammelt,
geordnet und für alle zugänglich gemacht wird. Den Ausweis habe ich
noch. Diese frühe Nähe zu Büchern und Bibliotheken hat vermutlich meinen
weiteren Weg mit beeinflusst. Im Studium habe ich mich zunächst dem
Archivbereich zugewandt und dort auch auf Anregung meines Vorgesetzten
praktische Erfahrungen gesammelt. Gleichzeitig ist mir immer klarer
geworden, dass ich genau an der Schnittstelle zwischen Wissen, Struktur
und Menschen arbeiten möchte. Das Gefühl, Wissen zu entdecken und
weiterzugeben, begleitet mich bis heute. Schließlich hat mich mein Weg
unter anderem nach Köln geführt, wo ich meine fachliche Ausrichtung im
Bibliotheks- und Informationsstudium vertiefen konnte, und von dort,
über meine weiteren beruflichen Stationen, bis heute zur DNB.

\begin{figure}
\centering
\includegraphics[width=0.75\textwidth]{img/abb1.jpg}
\caption{Kathrin Brannemann, Foto: DNB, Stephan Jockel, CC-BY-SA 3.0 DE}
\alt{Porträtfoto von Kathrin Brannemann mit hellblauer Bluse, dunkelblauem Jacket, Brille und blonden, mittellangen Haaren.}
\end{figure}

\textbf{Mögen Sie unseren Leserinnen und Lesern kurz erzählen, welche
Stationen auf Ihrem Weg zur Direktorin unseres Frankfurter Hauses
prägend waren?}

Auf meinem Weg zur DNB gab es viele prägende Impulse. Ein zentraler
Motor war dabei immer meine Neugier -- sie hat mich dazu gebracht,
unterschiedliche Bereiche kennenzulernen und neue Perspektiven
einzunehmen. Ich habe in sehr verschiedenen Kontexten gearbeitet, unter
anderem in Archiven, bei einem Bibliotheksverbund, in einer
Spezialbibliothek, im universitären Bereich sowie in der
außeruniversitären Forschung. Diese Vielfalt hat meinen Blick auf
Bibliotheken und ihre Rollen stark geschärft.

Für mich als Führungskraft war meine Zeit an der Universität Göttingen
besonders prägend. Dort stand ich vor der Aufgabe, ein sehr großes Haus
mit knapp fünfhundert Mitarbeiter*innen, an dem zuvor ein Gremium von
drei Direktoren tätig gewesen war, zeitweise allein zu führen. In einer
Phase knapper werdender Ressourcen ging es darum, Orientierung zu geben,
Prioritäten zu setzen, Entscheidungen transparent zu machen und
Verantwortung zu übernehmen. Gleichzeitig habe ich gelernt, wie wichtig
und unerlässlich der Verlass auf die Kompetenz und das Engagement der
Mitarbeitenden ist. Mir wurde viel Vertrauen entgegengebracht, das ich
auch hoffentlich in geeigneter Form zurückgeben konnte. Diese Erfahrung
hat mein Verständnis von Führung, gemeinsamer Zusammenarbeit und
strategischer Ausrichtung nachhaltig geprägt.

\textbf{Wenn Sie einmal an diese ganze Zeit zurückdenken, was war das
schönste (das lustigste, das emotionalste) Erlebnis in Ihrem bisherigen
Bibliotheksleben?~}

Ein besonders schönes und zugleich emotionales Erlebnis war für mich ein
Besuch der internationalen Forschungsstation in Ny-Ålesund in der Arktis
im Jahr 2019. Als Leiterin der Bibliothek des Alfred-Wegener-Instituts
Helmholtz-Zentrum für Polar- und Meeresforschung (AWI) war ich damals
Teil einer Delegation und durfte zehn Tage lang die Forschungsarbeit
direkt vor Ort erleben -- an einem der abgelegensten Orte der Welt.

Die Bibliothek in Ny-Ålesund funktioniert ähnlich wie die Unterstützung
auf dem Forschungsschiff Polarstern des AWI: Sie stellt sicher, dass
Forschende aus vielen Nationen jederzeit Zugang zu verlässlicher
Information haben -- direkt auf dem Schreibtisch, selbst am
\enquote{Ende der Welt}.

In dieser Umgebung wurde für mich spürbar, dass Bibliothekarin zu sein
nichts Abstraktes ist. Unsere Arbeit hat eine direkte, konkrete Wirkung
auf wissenschaftliche Arbeit und gesellschaftlich hochrelevante Themen.
Persönlich habe ich in dieser Zeit unglaublich viel über den Klimawandel
gelernt, zahlreiche fachliche Vorträge gehört und erlebt, wie
international, kooperativ und engagiert Forschung dort oben
funktioniert. Dieses Erlebnis hat mich fachlich wie emotional nachhaltig
geprägt.

\textbf{Worin sehen Sie aktuell die größten Herausforderungen und
wichtigsten Entwicklungen für Bibliotheken und insbesondere für
Landesbibliotheken und vor allem die Deutsche Nationalbibliothek?}

Für die Beantwortung dieser Frage benötigt man eigentlich mehr Zeit, da
das Thema sehr facettenreich ist. Wenn ich mich auf eine Herausforderung
fokussieren darf, dann liegt eine der größeren Herausforderungen für
Bibliotheken seit Jahren in der konsequenten Bewältigung der digitalen
Transformation. Für Landesbibliotheken und insbesondere für die Deutsche
Nationalbibliothek bedeutet das vor allem, mit der stetig wachsenden
Menge digitaler Publikationen umzugehen, diese langfristig zu sichern
und dauerhaft zugänglich zu machen. Die Sammlung, Erschließung und
Archivierung von Netzpublikationen, Forschungsdaten und digitalen
Kulturgütern stellt hohe technische, organisatorische und rechtliche
Anforderungen. Für die DNB steht dabei die Umsetzung ihrer strategischen
Kernaufgaben im Mittelpunkt: das gesamte publizierte, kulturelle Erbe
Deutschlands zu sammeln, dauerhaft zu bewahren, zu verzeichnen und
zugänglich zu machen -- unabhängig vom Format. Das Pflichtexemplarrecht
bildet dafür nach wie vor die unverzichtbare Grundlage, muss aber
kontinuierlich an digitalen Publikationsformen weiterentwickelt und
praktisch wirksam umgesetzt werden. Gleichzeitig geht es darum, digitale
Langzeitarchivierung, offene und verlässliche Zugänge, internationale
Vernetzung sowie eine klare Nutzerorientierung zusammenzudenken. So
sollten Daten aktiv strukturiert, vernetzt und für unterschiedliche
Zielgruppen nutzbar gemacht werden. Offene Schnittstellen,
standardisierte Metadaten und die stärkere Einbindung in nationale und
sogar internationale Dateninfrastrukturen können dies unterstützen. Die
Gemeinsame Normdatei (GND) ist zum Beispiel ein Dienst, um Normdaten
kooperativ nutzen und verwalten zu können. All dies erfordert
Priorisierung, Kooperation und ein hohes Maß an strategischer Klarheit
wie sie in den Strategischen Prioritäten 2025-2027 der DNB angelegt
sind. Nicht zuletzt ist auch der Kompetenzwandel eine wichtige
Herausforderung: Die strategische Ausrichtung der DNB betont die
Bedeutung interdisziplinärer Zusammenarbeit und neuer Qualifikationen an
der Schnittstelle von Bibliothekswesen, IT und Datenmanagement.

\textbf{Ganz offen gedacht: Was sollte aus Ihrer Sicht unbedingt bewahrt
werden?}

Ganz offen gedacht sollte alles bewahrt werden, was unsere Kultur, unser
Wissen und unsere Geschichte greifbar macht -- unabhängig davon, ob es
gedruckt, digital, bildlich oder auditiv ist. Das reicht von Büchern,
Zeitschriften und wissenschaftlichen Publikationen, Musik über
audiovisuelle Medien bis hin zu digitalen Inhalten, die heute alltäglich
sind, morgen aber verloren gegangen sein könnten. Besonders wichtig ist
mir, dass wir auch die Vielfalt der Perspektiven und Erfahrungen
dokumentieren: nicht nur das Offizielle, sondern auch das, was oft am
Rande steht -- lokale Geschichte, kulturelle Nischen oder Forschung, die
noch nicht Mainstream ist. Denn Bibliotheken und andere
Gedächtnisinstitutionen sind nicht nur Speicher, sondern Spiegel unserer
Gesellschaft und unserer Zeit und müssen den gesellschaftlichen Dialog
suchen und fördern. Die DNB hat mit dem Deutschen Buch- und
Schriftmuseum, dem Deutschen Exilarchiv 1933--1945 und dem Deutschen
Musikarchiv organisatorische Einheiten im Haus, die diesen Prozess mit
unterstützen.

\textbf{Wo liegen Ihre Prioritäten in den ersten Monaten Ihrer neuen
Tätigkeit?}

In den ersten Monaten liegt mein Fokus vor allem darauf, zuzuhören,
kennenzulernen und zu verstehen -- die Häuser in Leipzig und Frankfurt,
die Mitarbeitenden, die bestehenden Strukturen und Arbeitsweisen. Eine
Faustregel, die ich aus Erfahrung mitbringe, ist: Die ersten sechs
Monate dienen dem Ankommen, nach etwa einem Jahr fühlt man sich dann
meistens in allen Belangen einer Einrichtung wirklich sattelfest.
Parallel möchte ich von Beginn an natürlich versuchen Impulse zu setzen,
Dinge zu bewegen und die strategische Weiterentwicklung aktiv mit zu
unterstützen, ohne überstürzt zu handeln. Es geht darum, ein gutes
Gleichgewicht zwischen Verstehen, Beziehungen und Vertrauen aufzubauen
und ins strategische Handeln zu finden.

\textbf{Gibt es einen Wert oder eine Überzeugung, die Sie in Ihrer
Arbeit besonders prägt?}

Ein Wert, der meine Arbeit besonders prägt, ist das Prinzip des
Verständnisses: Wo man Themen, Inhalte und Zusammenhänge nachvollziehen
kann, entsteht hoffentlich ein Verständnis -- und darauf lassen sich
fundierte Entscheidungen aufbauen. Für mich bedeutet das auch
Transparenz: offen zu kommunizieren, was machbar ist, was aber auch
nicht und warum, um möglichst die Mitarbeitenden und externe
Partner*innen mitzunehmen. Natürlich gelingt das nicht immer perfekt,
aber es ist ein Leitprinzip, das meine Arbeitsweise prägt und den
Austausch, die Zusammenarbeit und die Akzeptanz von Entscheidungen
unterstützt. Kenntnis schafft Verständnis.

\textbf{Sinnvolles Schaffen -- was bedeutet das für Sie, wenn Sie an
Ihre neue Tätigkeit in der Deutschen Nationalbibliothek denken?}

Sinnvolles Schaffen bedeutet für mich nach Außen gerichtet, an etwas zu
arbeiten, das über den unmittelbaren Moment hinaus Wirkung entfaltet. In
der Deutschen Nationalbibliothek sehe ich genau diese Sinnhaftigkeit:
Sie bewahrt das kulturelle und geistige Erbe unseres Landes und macht es
für heutige und zukünftige Generationen zugänglich. Auf diesem Wege
fördert und unterstützt sie die Informations- und Meinungsfreiheit in
diesem Land. ‚Sinnvoll' ist für mich also eng verbunden mit Relevanz,
Nachhaltigkeit und der spürbaren Wirkung unserer Arbeit.

Zudem bedeutet sinnvolles Schaffen für mich auch, Verantwortung zu
übernehmen: für Qualität, für Kooperationen im nationalen und
internationalen Kontext und für die kontinuierliche Weiterentwicklung
der Bibliothek als moderne Gedächtnisinstitution. In meiner neuen
Tätigkeit möchte ich dazu beitragen, die strategischen Ziele der DNB mit
Leben zu füllen und gemeinsam Lösungen zu entwickeln, die Bestand haben.
Nach innen gerichtet bedeutet Sinnvolles schaffen für mich dabei, die
Mitarbeitenden in diese Prozesse einzubeziehen, dass sie ihr Know-How
und Kompetenzen einbringen können und ihre Fähigkeiten zu stärken.

%autor
\begin{center}\rule{0.5\linewidth}{0.5pt}\end{center}

\textbf{Kathrin Brannemann} ist
seit 1. Februar 2026 Direktorin der Deutschen Nationalbibliothek in
Frankfurt am Main.

\textbf{Josephine Kreutzer}, Referatsleitung Kommunikation, Deutsche
Nationalbibliothek.

\end{document}